% AML (CSAI2017P) --- Theory Quiz 3 (T03) --- 8 sets, A--H, four pages each.
% Generated from the item store; do not edit this file directly.
% Do not release until every group has sat it. Worked solutions: T03-Solutions.pdf.
\documentclass[10pt,a4paper]{article}
\usepackage[margin=16mm,top=12mm,bottom=14mm]{geometry}
\usepackage[T1]{fontenc}
\usepackage{lmodern}
\usepackage{enumitem}
\usepackage{fancyvrb}
\usepackage{amsmath}
\usepackage{amssymb}
\usepackage{array}
\usepackage{fancyhdr}
\usepackage{microtype}

\setlength{\parindent}{0pt}
\setlength{\emergencystretch}{3em}
\newcommand{\BL}{\rule[-1.5pt]{22mm}{0.4pt}}
\newcommand{\blk}[1]{\rule[-1.5pt]{#1}{0.4pt}}

\pagestyle{fancy}\fancyhf{}
\renewcommand{\headrulewidth}{0.4pt}
\renewcommand{\footrulewidth}{0.4pt}
\newcommand{\SetLetter}{}
\fancyhead[L]{\footnotesize CSAI2017P --- Applied Machine Learning (Theory) --- Theory Quiz 3}
\fancyhead[R]{\footnotesize\bfseries Set \SetLetter}
\fancyfoot[L]{\footnotesize Name:\,\blk{34mm} \quad SAP ID:\,\blk{28mm}}
\fancyfoot[R]{\footnotesize Set \SetLetter\ --- page \thepage\ of 4}

\DefineVerbatimEnvironment{code}{Verbatim}%
  {fontsize=\small,xleftmargin=5mm,baselinestretch=0.98}

\newlist{qA}{enumerate}{1}
\setlist[qA]{label=\textbf{A\arabic*.},leftmargin=8mm,topsep=4pt,itemsep=9pt,parsep=1pt}
\newlist{qB}{enumerate}{1}
\setlist[qB]{label=\textbf{B\arabic*.},leftmargin=8mm,topsep=5pt,itemsep=13pt,parsep=1pt}
\newlist{qC}{enumerate}{1}
\setlist[qC]{label=\textbf{C\arabic*.},leftmargin=8mm,topsep=5pt,itemsep=8pt,parsep=1pt}
\newlist{qD}{enumerate}{1}
\setlist[qD]{label=\textbf{D\arabic*.},leftmargin=8mm,topsep=3pt,itemsep=4pt,parsep=1pt}
\newlist{qE}{enumerate}{1}
\setlist[qE]{label=\textbf{E\arabic*.},leftmargin=8mm,topsep=3pt,itemsep=5pt,parsep=1pt}

\newcommand{\parthead}[2]{\vspace{2pt}\noindent\rule{\linewidth}{0.4pt}\par\vspace{2pt}%
  \noindent{\bfseries Part #1}\hfill{\small #2}\par\vspace{1pt}}

\makeatletter
\newcommand{\rules}[2]{\par\vspace{2mm}%
  \@tempcnta=0
  \loop\ifnum\@tempcnta<#1\relax
    \noindent\rule{\linewidth}{0.3pt}\par\vspace{#2}%
    \advance\@tempcnta by 1
  \repeat}
\makeatother

\begin{document}


\renewcommand{\SetLetter}{A}\setcounter{page}{1}


\begin{center}
{\large\bfseries CSAI2017P --- Applied Machine Learning (Theory)}\\[2pt]
{\large\bfseries Theory Quiz 3 (T03) \quad\textbullet\quad Set A}\\[3pt]
{\small Duration \textbf{40 minutes} \quad\textbullet\quad Maximum marks \textbf{20}
 \quad\textbullet\quad 15 questions on four pages}
\end{center}
\vspace{2pt}

% Column widths differ slightly from LQ01's: at LQ01's widths the row runs
% ~11.6pt past the margin and "Signature" does not fit its own cell.
\noindent\begin{tabular}{@{}p{1.4cm}p{5.0cm}p{1.9cm}p{3.5cm}p{1.3cm}p{2.4cm}@{}}
\textbf{Name} & \hrulefill & \textbf{SAP ID} & \hrulefill & \textbf{Batch} & \hrulefill\\[9pt]
\textbf{Date} & \hrulefill & \textbf{Signature} & \hrulefill & \textbf{Set} & \fbox{\rule{0pt}{11pt}\ A\ }\\
\end{tabular}

\vspace{4pt}
\noindent\fbox{\parbox{\dimexpr\linewidth-2\fboxsep-2\fboxrule}{\small
\textbf{Syllabus.} Prerequisite material --- probability, Python and elementary calculus --- together with
Units I--III (Lectures 1--3) as retained material: generalisation and data leakage, loss against metric, and
batch sizes and momentum; and, principally, \textbf{Unit III (Lecture 4)}: adaptive optimisers --- Adagrad,
RMSprop, Adadelta and Adam, bias correction, and what adaptivity does and does not buy.\\[3pt]
\textbf{Instructions.} Closed book. Write your answers on this paper, in the space provided. A
calculator is not needed --- the arithmetic is deliberately short. Answer every question --- there
is no negative marking. In Part C a reason is \textbf{optional}, but a correct reason can rescue a
wrong answer. In Parts D and E the marks are for the reasoning, not the vocabulary; an answer that
states what it cannot determine earns credit, one that asserts without a reason does not.\\[3pt]
\textbf{About the set letter.} Sets are handed out \textbf{at random}. Every set covers the same
topics at the same level, for the same marks, against the same marking scheme.}}
\vspace{3pt}


\parthead{A --- Multiple choice}{5 questions $\times$ 1 = 5 marks. Circle one letter.}

\begin{qA}
  \item One card is drawn at random from a standard $52$-card pack. Let $H$ be the event that it is a heart. Then $P(H)$ equals\\
    (a) $1/13$ \quad (b) $1/2$\\
    (c) $1/4$ \quad (d) $1/52$
  \item \texttt{len(range(1, 20, 4))} evaluates to\\
    (a) 5 \quad (b) 4\\
    (c) 6 \quad (d) an error
  \item With $n$ training examples and batch size $B$, the number of parameter updates in one epoch is\\
    (a) $n$ \quad (b) $B$\\
    (c) $1$, whatever $B$ is \quad (d) $n/B$
  \item A step that \emph{learns} from the data is fitted on the whole of $X$ before $X$ is split. The consequence is\\
    (a) the model underfits \quad (b) information from the test rows influences training\\
    (c) the optimiser converges more slowly \quad (d) nothing --- the split is still random
  \item In Adagrad, $s_t = s_{t-1} + g_t \odot g_t$. The quantity $s_t$ is therefore\\
    (a) a running \emph{average} of squared gradients \quad (b) the running sum of the gradients themselves\\
    (c) the running \emph{sum} of squared gradients \quad (d) the current gradient, squared
\end{qA}

\vspace{4pt}

\parthead{B --- Fill in the blanks}{4 $\times$ 1 = 4 marks. Write the missing word on the rule.}

\begin{qB}

  \item The derivative of $w^3$ with respect to $w$ is \BL.

  \item Averaging the gradient over a mini-batch of size $B$ reduces its noise by a factor of \BL.

  \item Adagrad's accumulator never decreases, so its effective step size decays like $1/\BL$.

  \item Because $m$ and $v$ are corrected for their start at zero, Adam's first step has size exactly \BL.

\end{qB}

\vfill
\begin{center}\small
\begin{tabular}{|>{\centering\arraybackslash}m{1.9cm}|*{5}{>{\centering\arraybackslash}m{1.6cm}|}}
\hline
\rule{0pt}{12pt}\textbf{Question} & A1 & A2 & A3 & A4 & A5\\[2pt]\hline
\rule{0pt}{20pt}\textbf{Answer} & & & & & \\[7pt]\hline
\end{tabular}\\[2pt]
{\footnotesize Copy your circled Part A letters into this grid. The grid is what is marked.}
\end{center}

\newpage

\parthead{C --- True or false}{3 $\times$ 1 = 3 marks. Circle one. A reason is optional, but write one --- it can rescue a wrong answer.}

\begin{qC}

  \item Data leakage announces itself through a \emph{poor} score on the held-out data. \hfill \textbf{TRUE} \; / \; \textbf{FALSE}\\[1pt]
    {\small Reason:}\,\blk{140mm}\rules{1}{6.4mm}

  \item With the learning rate tuned separately for each, an adaptive optimiser will beat plain gradient descent with momentum on any problem. \hfill \textbf{TRUE} \; / \; \textbf{FALSE}\\[1pt]
    {\small Reason:}\,\blk{140mm}\rules{1}{6.4mm}

  \item On a ravine-shaped loss surface, plain momentum can converge in fewer iterations than Adam. \hfill \textbf{TRUE} \; / \; \textbf{FALSE}\\[1pt]
    {\small Reason:}\,\blk{140mm}\rules{1}{6.4mm}

\end{qC}

\vspace{4pt}

\parthead{D --- Short answer}{2 questions $\times$ 2 = 4 marks. Show how you decide, not just what you decide.}

\begin{qD}

  \item Both optimisers start from $s_0 = 0$ and take their \emph{first} step with $\eta = 0.5$ and gradient $g_1 = -6$. Adagrad uses $s_t = s_{t-1} + g_t^2$; RMSprop uses $s_t = \gamma s_{t-1} + (1-\gamma) g_t^2$ with $\gamma = 0.9$. Both then apply $w \leftarrow w - \eta g_t / \sqrt{s_t}$; take $\varepsilon$ as negligible.
    \begin{enumerate}[label=(\alph*),leftmargin=7mm,itemsep=1pt,topsep=2pt]
      \item Compute the size of each optimiser's first step, $\lvert \Delta w_1 \rvert$. Give two decimal places where the answer is not exact.
      \item Neither answer depends on how large $g_1$ was. Show why.
    \end{enumerate}\par\vspace{64mm}

\end{qD}

\newpage

\parthead{D --- Short answer, continued}{}

\begin{qD}[start=2]

  \item A classifier is trained by minimising cross-entropy, but the figure reported to the client is accuracy. \textbf{(a)} Explain why these are two different quantities rather than two names for one. \textbf{(b)} State which of them you use to decide the model is finished, and why.\par\vspace{80mm}

\end{qD}

\vspace{4pt}

\parthead{E --- Reasoning}{1 question $\times$ 4 = 4 marks. You are marked on the reasoning you make visible.}

\begin{qE}

  \item A classmate asserts: \emph{``Adam always beats plain gradient descent with momentum.''} You have one
afternoon and a laptop. Design the smallest experiment that would actually settle it.


    \vspace{2mm}
    Answer all four parts overleaf. There is no single right experiment --- you are marked on whether
    the one you describe could actually decide the question.
    \begin{enumerate}[label=(\roman*),leftmargin=7mm,itemsep=2pt,topsep=2pt]
      \item What would you run it on? Say how many problems and of what kind. \hfill \textbf{[1]}
      \item What would you measure, and at what point? \hfill \textbf{[1]}
      \item Name the one thing that must be held equal between the two optimisers, without which the
            comparison means nothing at all. \hfill \textbf{[1]}
      \item State the result that would make you accept the claim, \emph{and} the result that would
            make you reject it. \hfill \textbf{[1]}
    \end{enumerate}

\end{qE}

\newpage

\parthead{E --- Reasoning, answer space}{Label each part (i), (ii), (iii). A sketch, a table or a list is as acceptable as prose.}

\vfill

{\footnotesize\itshape Continue on a continuation sheet if you need more room; label anything you write there with its question number.}

\newpage

\renewcommand{\SetLetter}{B}\setcounter{page}{1}


\begin{center}
{\large\bfseries CSAI2017P --- Applied Machine Learning (Theory)}\\[2pt]
{\large\bfseries Theory Quiz 3 (T03) \quad\textbullet\quad Set B}\\[3pt]
{\small Duration \textbf{40 minutes} \quad\textbullet\quad Maximum marks \textbf{20}
 \quad\textbullet\quad 15 questions on four pages}
\end{center}
\vspace{2pt}

% Column widths differ slightly from LQ01's: at LQ01's widths the row runs
% ~11.6pt past the margin and "Signature" does not fit its own cell.
\noindent\begin{tabular}{@{}p{1.4cm}p{5.0cm}p{1.9cm}p{3.5cm}p{1.3cm}p{2.4cm}@{}}
\textbf{Name} & \hrulefill & \textbf{SAP ID} & \hrulefill & \textbf{Batch} & \hrulefill\\[9pt]
\textbf{Date} & \hrulefill & \textbf{Signature} & \hrulefill & \textbf{Set} & \fbox{\rule{0pt}{11pt}\ B\ }\\
\end{tabular}

\vspace{4pt}
\noindent\fbox{\parbox{\dimexpr\linewidth-2\fboxsep-2\fboxrule}{\small
\textbf{Syllabus.} Prerequisite material --- probability, Python and elementary calculus --- together with
Units I--III (Lectures 1--3) as retained material: generalisation and data leakage, loss against metric, and
batch sizes and momentum; and, principally, \textbf{Unit III (Lecture 4)}: adaptive optimisers --- Adagrad,
RMSprop, Adadelta and Adam, bias correction, and what adaptivity does and does not buy.\\[3pt]
\textbf{Instructions.} Closed book. Write your answers on this paper, in the space provided. A
calculator is not needed --- the arithmetic is deliberately short. Answer every question --- there
is no negative marking. In Part C a reason is \textbf{optional}, but a correct reason can rescue a
wrong answer. In Parts D and E the marks are for the reasoning, not the vocabulary; an answer that
states what it cannot determine earns credit, one that asserts without a reason does not.\\[3pt]
\textbf{About the set letter.} Sets are handed out \textbf{at random}. Every set covers the same
topics at the same level, for the same marks, against the same marking scheme.}}
\vspace{3pt}


\parthead{A --- Multiple choice}{5 questions $\times$ 1 = 5 marks. Circle one letter.}

\begin{qA}
  \item A crate holds $20$ items, of which $5$ are marked for rework. One is drawn at random. The probability that it is marked is\\
    (a) $1/5$ \quad (b) $1/20$\\
    (c) $3/4$ \quad (d) $1/4$
  \item For \texttt{xs = [2, 4, 6, 8, 10]}, the expression \texttt{len(xs[1:4])} evaluates to\\
    (a) 2 \quad (b) 3\\
    (c) 4 \quad (d) 5
  \item A table records $n = 4000$ rows and $B = 32$. The number of parameter updates in one epoch is\\
    (a) 125 \quad (b) 4000\\
    (c) 32 \quad (d) 1
  \item A table shows a scaler fitted on all $500$ rows, then a $400/100$ split, then a test score of $0.99$. The most likely reading of that $0.99$ is\\
    (a) the model is genuinely excellent \quad (b) the model has underfitted\\
    (c) it is optimistic --- test information reached the model \quad (d) the test set is too large
  \item Two accumulators are tabulated at each step: one column only ever grows, the other rises and falls. The column that only grows belongs to\\
    (a) RMSprop \quad (b) Adadelta\\
    (c) Adam \quad (d) Adagrad
\end{qA}

\vspace{4pt}

\parthead{B --- Fill in the blanks}{4 $\times$ 1 = 4 marks. Write the missing word on the rule.}

\begin{qB}

  \item A table of $w$ against $w^3$ has gradient column \BL.

  \item Measured on the diabetes data, the RMS gradient error fell from $6.32$ at $B=1$ to $1.06$ at $B=32$ --- a factor of $5.9$, because the noise falls like $1/\BL$.

  \item A column of Adagrad's effective step size falls by $9.4\times$ over $2800$ updates and keeps falling, because the step decays like $1/\BL$.

  \item A table of Adam's first few steps reads $0.1000$, $0.0999$, $0.0997$ at $\eta = 0.1$: the corrected first step has size exactly \BL.

\end{qB}

\vfill
\begin{center}\small
\begin{tabular}{|>{\centering\arraybackslash}m{1.9cm}|*{5}{>{\centering\arraybackslash}m{1.6cm}|}}
\hline
\rule{0pt}{12pt}\textbf{Question} & A1 & A2 & A3 & A4 & A5\\[2pt]\hline
\rule{0pt}{20pt}\textbf{Answer} & & & & & \\[7pt]\hline
\end{tabular}\\[2pt]
{\footnotesize Copy your circled Part A letters into this grid. The grid is what is marked.}
\end{center}

\newpage

\parthead{C --- True or false}{3 $\times$ 1 = 3 marks. Circle one. A reason is optional, but write one --- it can rescue a wrong answer.}

\begin{qC}

  \item A results table in which the held-out score is \emph{worse} than expected is the signature of data leakage. \hfill \textbf{TRUE} \; / \; \textbf{FALSE}\\[1pt]
    {\small Reason:}\,\blk{140mm}\rules{1}{6.4mm}

  \item A table comparing six optimisers, each at its own tuned learning rate, must show an adaptive method at the top. \hfill \textbf{TRUE} \; / \; \textbf{FALSE}\\[1pt]
    {\small Reason:}\,\blk{140mm}\rules{1}{6.4mm}

  \item A table reading $135$ iterations for momentum against $262$ for Adam on a ravine is possible. \hfill \textbf{TRUE} \; / \; \textbf{FALSE}\\[1pt]
    {\small Reason:}\,\blk{140mm}\rules{1}{6.4mm}

\end{qC}

\vspace{4pt}

\parthead{D --- Short answer}{2 questions $\times$ 2 = 4 marks. Show how you decide, not just what you decide.}

\begin{qD}

  \item The table below is to be completed for the \emph{first} step of each optimiser from $s_0 = 0$, with $\eta = 0.2$, $\gamma = 0.75$ and $g_1 = +4$. Adagrad accumulates $s_t = s_{t-1} + g_t^2$; RMSprop accumulates $s_t = \gamma s_{t-1} + (1-\gamma)g_t^2$; both then apply $w \leftarrow w - \eta g_t/\sqrt{s_t}$, with $\varepsilon$ negligible.\par\vspace{1mm}\hspace*{6mm}\begin{tabular}{@{}lll@{}}\hline\rule{0pt}{10pt}Optimiser & $s_1$ & $\lvert \Delta w_1 \rvert$\\[1pt]\hline\rule{0pt}{10pt}Adagrad & & \\ RMSprop & & \\[1pt]\hline\end{tabular}
    \begin{enumerate}[label=(\alph*),leftmargin=7mm,itemsep=1pt,topsep=2pt]
      \item Compute the size of each optimiser's first step, $\lvert \Delta w_1 \rvert$. Give two decimal places where the answer is not exact.
      \item Neither answer depends on how large $g_1$ was. Show why.
    \end{enumerate}\par\vspace{64mm}

\end{qD}

\newpage

\parthead{D --- Short answer, continued}{}

\begin{qD}[start=2]

  \item A results table carries two columns for the same classifier: a cross-entropy column, which the training loop minimised, and an accuracy column, which is what the client was shown. \textbf{(a)} Explain why these are two different quantities rather than two names for one. \textbf{(b)} State which of them you use to decide the model is finished, and why.\par\vspace{80mm}

\end{qD}

\vspace{4pt}

\parthead{E --- Reasoning}{1 question $\times$ 4 = 4 marks. You are marked on the reasoning you make visible.}

\begin{qE}

  \item A results table circulating in the lab is being used to argue that \emph{``Adam always beats plain
gradient descent with momentum''}. You have one afternoon and a laptop. Design the smallest experiment that
would actually settle it.


    \vspace{2mm}
    Answer all four parts overleaf. There is no single right experiment --- you are marked on whether
    the one you describe could actually decide the question.
    \begin{enumerate}[label=(\roman*),leftmargin=7mm,itemsep=2pt,topsep=2pt]
      \item What would you run it on? Say how many problems and of what kind. \hfill \textbf{[1]}
      \item What would you measure, and at what point? \hfill \textbf{[1]}
      \item Name the one thing that must be held equal between the two optimisers, without which the
            comparison means nothing at all. \hfill \textbf{[1]}
      \item State the result that would make you accept the claim, \emph{and} the result that would
            make you reject it. \hfill \textbf{[1]}
    \end{enumerate}

\end{qE}

\newpage

\parthead{E --- Reasoning, answer space}{Label each part (i), (ii), (iii). A sketch, a table or a list is as acceptable as prose.}

\vfill

{\footnotesize\itshape Continue on a continuation sheet if you need more room; label anything you write there with its question number.}

\newpage

\renewcommand{\SetLetter}{C}\setcounter{page}{1}


\begin{center}
{\large\bfseries CSAI2017P --- Applied Machine Learning (Theory)}\\[2pt]
{\large\bfseries Theory Quiz 3 (T03) \quad\textbullet\quad Set C}\\[3pt]
{\small Duration \textbf{40 minutes} \quad\textbullet\quad Maximum marks \textbf{20}
 \quad\textbullet\quad 15 questions on four pages}
\end{center}
\vspace{2pt}

% Column widths differ slightly from LQ01's: at LQ01's widths the row runs
% ~11.6pt past the margin and "Signature" does not fit its own cell.
\noindent\begin{tabular}{@{}p{1.4cm}p{5.0cm}p{1.9cm}p{3.5cm}p{1.3cm}p{2.4cm}@{}}
\textbf{Name} & \hrulefill & \textbf{SAP ID} & \hrulefill & \textbf{Batch} & \hrulefill\\[9pt]
\textbf{Date} & \hrulefill & \textbf{Signature} & \hrulefill & \textbf{Set} & \fbox{\rule{0pt}{11pt}\ C\ }\\
\end{tabular}

\vspace{4pt}
\noindent\fbox{\parbox{\dimexpr\linewidth-2\fboxsep-2\fboxrule}{\small
\textbf{Syllabus.} Prerequisite material --- probability, Python and elementary calculus --- together with
Units I--III (Lectures 1--3) as retained material: generalisation and data leakage, loss against metric, and
batch sizes and momentum; and, principally, \textbf{Unit III (Lecture 4)}: adaptive optimisers --- Adagrad,
RMSprop, Adadelta and Adam, bias correction, and what adaptivity does and does not buy.\\[3pt]
\textbf{Instructions.} Closed book. Write your answers on this paper, in the space provided. A
calculator is not needed --- the arithmetic is deliberately short. Answer every question --- there
is no negative marking. In Part C a reason is \textbf{optional}, but a correct reason can rescue a
wrong answer. In Parts D and E the marks are for the reasoning, not the vocabulary; an answer that
states what it cannot determine earns credit, one that asserts without a reason does not.\\[3pt]
\textbf{About the set letter.} Sets are handed out \textbf{at random}. Every set covers the same
topics at the same level, for the same marks, against the same marking scheme.}}
\vspace{3pt}


\parthead{A --- Multiple choice}{5 questions $\times$ 1 = 5 marks. Circle one letter.}

\begin{qA}
  \item A spinner is divided into eight equal sectors, two of them shaded. The probability that one spin lands on a shaded sector is\\
    (a) $1/4$ \quad (b) $1/8$\\
    (c) $1/2$ \quad (d) $3/4$
  \item \texttt{len(set('banana'))} evaluates to\\
    (a) 6 \quad (b) 2\\
    (c) 3 \quad (d) an error
  \item A bar chart shows the number of weight updates in one epoch, for three batch sizes run on the same $4000$-row dataset. The bar for $B = 32$ has height\\
    (a) 4000 \quad (b) 125\\
    (c) 32 \quad (d) 1
  \item A learning curve rises smoothly to a test score far above anything the domain expert expects. Before believing it, the first thing to suspect is\\
    (a) too small a learning rate \quad (b) too large a batch size\\
    (c) an unlucky random seed \quad (d) a fitted step that saw the whole dataset before the split
  \item A plot of the denominator $\sqrt{s_t}$ against $t$ is monotonically increasing and never comes back down. The optimiser is\\
    (a) Adagrad \quad (b) RMSprop\\
    (c) Adadelta \quad (d) Adam
\end{qA}

\vspace{4pt}

\parthead{B --- Fill in the blanks}{4 $\times$ 1 = 4 marks. Write the missing word on the rule.}

\begin{qB}

  \item The slope of the curve $y = w^3$ at a point $w$ is \BL.

  \item A plot of gradient noise against batch size falls as a power law with exponent $-1/2$: the noise goes like $1/\BL$.

  \item A plot of Adagrad's effective step against $t$ is a straight line of slope $-1/2$ on log--log axes, because the step decays like $1/\BL$.

  \item A plot of Adam's step size against $t$ starts flat at the value $\eta$ rather than above it, because bias correction makes the first step exactly \BL.

\end{qB}

\vfill
\begin{center}\small
\begin{tabular}{|>{\centering\arraybackslash}m{1.9cm}|*{5}{>{\centering\arraybackslash}m{1.6cm}|}}
\hline
\rule{0pt}{12pt}\textbf{Question} & A1 & A2 & A3 & A4 & A5\\[2pt]\hline
\rule{0pt}{20pt}\textbf{Answer} & & & & & \\[7pt]\hline
\end{tabular}\\[2pt]
{\footnotesize Copy your circled Part A letters into this grid. The grid is what is marked.}
\end{center}

\newpage

\parthead{C --- True or false}{3 $\times$ 1 = 3 marks. Circle one. A reason is optional, but write one --- it can rescue a wrong answer.}

\begin{qC}

  \item On a plot of held-out score against epoch, data leakage shows up as a curve sitting \emph{below} where it should. \hfill \textbf{TRUE} \; / \; \textbf{FALSE}\\[1pt]
    {\small Reason:}\,\blk{140mm}\rules{1}{6.4mm}

  \item If six optimisers are each plotted at their own best learning rate, the adaptive ones must end up below the non-adaptive ones. \hfill \textbf{TRUE} \; / \; \textbf{FALSE}\\[1pt]
    {\small Reason:}\,\blk{140mm}\rules{1}{6.4mm}

  \item A trajectory plot in which momentum reaches the minimum of a narrow valley in fewer steps than Adam is possible. \hfill \textbf{TRUE} \; / \; \textbf{FALSE}\\[1pt]
    {\small Reason:}\,\blk{140mm}\rules{1}{6.4mm}

\end{qC}

\vspace{4pt}

\parthead{D --- Short answer}{2 questions $\times$ 2 = 4 marks. Show how you decide, not just what you decide.}

\begin{qD}

  \item Two optimisers are to be drawn taking their \emph{first} step from the same point, with $s_0 = 0$, $\eta = 0.4$ and a gradient of $g_1 = -10$ at that point. Adagrad accumulates $s_t = s_{t-1} + g_t^2$; RMSprop accumulates $s_t = \gamma s_{t-1} + (1-\gamma)g_t^2$ with $\gamma = 0.84$. Both then apply $w \leftarrow w - \eta g_t/\sqrt{s_t}$; take $\varepsilon$ as negligible.
    \begin{enumerate}[label=(\alph*),leftmargin=7mm,itemsep=1pt,topsep=2pt]
      \item Compute the size of each optimiser's first step, $\lvert \Delta w_1 \rvert$. Give two decimal places where the answer is not exact.
      \item Neither answer depends on how large $g_1$ was. Show why.
    \end{enumerate}\par\vspace{64mm}

\end{qD}

\newpage

\parthead{D --- Short answer, continued}{}

\begin{qD}[start=2]

  \item Two curves are plotted against the epoch: the cross-entropy the training loop was minimising, and the accuracy that was reported to the client. They do not have the same shape. \textbf{(a)} Explain why these are two different quantities rather than two names for one. \textbf{(b)} State which of them you use to decide the model is finished, and why.\par\vspace{80mm}

\end{qD}

\vspace{4pt}

\parthead{E --- Reasoning}{1 question $\times$ 4 = 4 marks. You are marked on the reasoning you make visible.}

\begin{qE}

  \item A plot is being passed around as proof that \emph{``Adam always beats plain gradient descent with
momentum''}. You have one afternoon and a laptop. Design the smallest experiment that would actually settle
it.


    \vspace{2mm}
    Answer all four parts overleaf. There is no single right experiment --- you are marked on whether
    the one you describe could actually decide the question.
    \begin{enumerate}[label=(\roman*),leftmargin=7mm,itemsep=2pt,topsep=2pt]
      \item What would you run it on? Say how many problems and of what kind. \hfill \textbf{[1]}
      \item What would you measure, and at what point? \hfill \textbf{[1]}
      \item Name the one thing that must be held equal between the two optimisers, without which the
            comparison means nothing at all. \hfill \textbf{[1]}
      \item State the result that would make you accept the claim, \emph{and} the result that would
            make you reject it. \hfill \textbf{[1]}
    \end{enumerate}

\end{qE}

\newpage

\parthead{E --- Reasoning, answer space}{Label each part (i), (ii), (iii). A sketch, a table or a list is as acceptable as prose.}

\vfill

{\footnotesize\itshape Continue on a continuation sheet if you need more room; label anything you write there with its question number.}

\newpage

\renewcommand{\SetLetter}{D}\setcounter{page}{1}


\begin{center}
{\large\bfseries CSAI2017P --- Applied Machine Learning (Theory)}\\[2pt]
{\large\bfseries Theory Quiz 3 (T03) \quad\textbullet\quad Set D}\\[3pt]
{\small Duration \textbf{40 minutes} \quad\textbullet\quad Maximum marks \textbf{20}
 \quad\textbullet\quad 15 questions on four pages}
\end{center}
\vspace{2pt}

% Column widths differ slightly from LQ01's: at LQ01's widths the row runs
% ~11.6pt past the margin and "Signature" does not fit its own cell.
\noindent\begin{tabular}{@{}p{1.4cm}p{5.0cm}p{1.9cm}p{3.5cm}p{1.3cm}p{2.4cm}@{}}
\textbf{Name} & \hrulefill & \textbf{SAP ID} & \hrulefill & \textbf{Batch} & \hrulefill\\[9pt]
\textbf{Date} & \hrulefill & \textbf{Signature} & \hrulefill & \textbf{Set} & \fbox{\rule{0pt}{11pt}\ D\ }\\
\end{tabular}

\vspace{4pt}
\noindent\fbox{\parbox{\dimexpr\linewidth-2\fboxsep-2\fboxrule}{\small
\textbf{Syllabus.} Prerequisite material --- probability, Python and elementary calculus --- together with
Units I--III (Lectures 1--3) as retained material: generalisation and data leakage, loss against metric, and
batch sizes and momentum; and, principally, \textbf{Unit III (Lecture 4)}: adaptive optimisers --- Adagrad,
RMSprop, Adadelta and Adam, bias correction, and what adaptivity does and does not buy.\\[3pt]
\textbf{Instructions.} Closed book. Write your answers on this paper, in the space provided. A
calculator is not needed --- the arithmetic is deliberately short. Answer every question --- there
is no negative marking. In Part C a reason is \textbf{optional}, but a correct reason can rescue a
wrong answer. In Parts D and E the marks are for the reasoning, not the vocabulary; an answer that
states what it cannot determine earns credit, one that asserts without a reason does not.\\[3pt]
\textbf{About the set letter.} Sets are handed out \textbf{at random}. Every set covers the same
topics at the same level, for the same marks, against the same marking scheme.}}
\vspace{3pt}


\parthead{A --- Multiple choice}{5 questions $\times$ 1 = 5 marks. Circle one letter.}

\begin{qA}
  \item A class of $40$ students includes $10$ who cycle to campus. One student is picked at random. The probability that this student cycles is\\
    (a) $1/10$ \quad (b) $1/4$\\
    (c) $1/40$ \quad (d) $3/4$
  \item A shop's price list is written as \texttt{\{'tea': 20, 'milk': 30, 'tea': 25\}}. Its \texttt{len} is\\
    (a) 3 \quad (b) 1\\
    (c) an error \quad (d) 2
  \item An engineer has $4000$ training rows and processes them $32$ at a time. In one pass through the data the model's weights are updated\\
    (a) 4000 times \quad (b) 32 times\\
    (c) 125 times \quad (d) once
  \item An intern reports a test accuracy far above what the team expected. Asked how, she says she scaled the whole table before splitting it. The reported figure is\\
    (a) optimistic --- the test rows helped fit the scaler \quad (b) trustworthy --- scaling is not a model\\
    (c) pessimistic --- scaling loses information \quad (d) unaffected --- scaling is reversible
  \item A colleague wants an optimiser that needs no learning rate at all, because the one he picks is always wrong. He wants\\
    (a) Adagrad \quad (b) Adadelta\\
    (c) RMSprop \quad (d) Adam
\end{qA}

\vspace{4pt}

\parthead{B --- Fill in the blanks}{4 $\times$ 1 = 4 marks. Write the missing word on the rule.}

\begin{qB}

  \item A cost grows as the cube of a setting $w$. Its rate of change with $w$ is \BL.

  \item An engineer who quadruples the batch size should expect the gradient noise to halve, because noise falls like $1/\BL$.

  \item A colleague complains that after a long run his model has almost stopped moving. He is using Adagrad, whose step decays like $1/\BL$.

  \item An engineer notes that whatever the first gradient happens to be, Adam's first move has size exactly \BL.

\end{qB}

\vfill
\begin{center}\small
\begin{tabular}{|>{\centering\arraybackslash}m{1.9cm}|*{5}{>{\centering\arraybackslash}m{1.6cm}|}}
\hline
\rule{0pt}{12pt}\textbf{Question} & A1 & A2 & A3 & A4 & A5\\[2pt]\hline
\rule{0pt}{20pt}\textbf{Answer} & & & & & \\[7pt]\hline
\end{tabular}\\[2pt]
{\footnotesize Copy your circled Part A letters into this grid. The grid is what is marked.}
\end{center}

\newpage

\parthead{C --- True or false}{3 $\times$ 1 = 3 marks. Circle one. A reason is optional, but write one --- it can rescue a wrong answer.}

\begin{qC}

  \item A team whose held-out score came out disappointingly low should suspect data leakage. \hfill \textbf{TRUE} \; / \; \textbf{FALSE}\\[1pt]
    {\small Reason:}\,\blk{140mm}\rules{1}{6.4mm}

  \item A manager who has tuned the learning rate for every optimiser on the shortlist can be confident the adaptive ones will win. \hfill \textbf{TRUE} \; / \; \textbf{FALSE}\\[1pt]
    {\small Reason:}\,\blk{140mm}\rules{1}{6.4mm}

  \item An engineer working on a long narrow valley may find plain momentum reaches the bottom in fewer iterations than Adam. \hfill \textbf{TRUE} \; / \; \textbf{FALSE}\\[1pt]
    {\small Reason:}\,\blk{140mm}\rules{1}{6.4mm}

\end{qC}

\vspace{4pt}

\parthead{D --- Short answer}{2 questions $\times$ 2 = 4 marks. Show how you decide, not just what you decide.}

\begin{qD}

  \item An engineer runs the same model twice from the same starting weights, once with Adagrad and once with RMSprop, and compares only the \emph{first} step. Both start from $s_0 = 0$ with $\eta = 0.25$, and the gradient at the starting point is $g_1 = +2$. Adagrad accumulates $s_t = s_{t-1} + g_t^2$; RMSprop accumulates $s_t = \gamma s_{t-1} + (1-\gamma)g_t^2$ with $\gamma = 0.9375$. Both apply $w \leftarrow w - \eta g_t/\sqrt{s_t}$; take $\varepsilon$ as negligible.
    \begin{enumerate}[label=(\alph*),leftmargin=7mm,itemsep=1pt,topsep=2pt]
      \item Compute the size of each optimiser's first step, $\lvert \Delta w_1 \rvert$. Give two decimal places where the answer is not exact.
      \item Neither answer depends on how large $g_1$ was. Show why.
    \end{enumerate}\par\vspace{64mm}

\end{qD}

\newpage

\parthead{D --- Short answer, continued}{}

\begin{qD}[start=2]

  \item A client has been shown an accuracy figure for a spam filter. The training loop that produced the filter never used accuracy at all --- it minimised cross-entropy. \textbf{(a)} Explain why these are two different quantities rather than two names for one. \textbf{(b)} State which of them you use to decide the model is finished, and why.\par\vspace{80mm}

\end{qD}

\vspace{4pt}

\parthead{E --- Reasoning}{1 question $\times$ 4 = 4 marks. You are marked on the reasoning you make visible.}

\begin{qE}

  \item A colleague insists, in a design meeting, that \emph{``Adam always beats plain gradient descent with
momentum''} and wants it written into the team's defaults. You have one afternoon and a laptop. Design the
smallest experiment that would actually settle it.


    \vspace{2mm}
    Answer all four parts overleaf. There is no single right experiment --- you are marked on whether
    the one you describe could actually decide the question.
    \begin{enumerate}[label=(\roman*),leftmargin=7mm,itemsep=2pt,topsep=2pt]
      \item What would you run it on? Say how many problems and of what kind. \hfill \textbf{[1]}
      \item What would you measure, and at what point? \hfill \textbf{[1]}
      \item Name the one thing that must be held equal between the two optimisers, without which the
            comparison means nothing at all. \hfill \textbf{[1]}
      \item State the result that would make you accept the claim, \emph{and} the result that would
            make you reject it. \hfill \textbf{[1]}
    \end{enumerate}

\end{qE}

\newpage

\parthead{E --- Reasoning, answer space}{Label each part (i), (ii), (iii). A sketch, a table or a list is as acceptable as prose.}

\vfill

{\footnotesize\itshape Continue on a continuation sheet if you need more room; label anything you write there with its question number.}

\newpage

\renewcommand{\SetLetter}{E}\setcounter{page}{1}


\begin{center}
{\large\bfseries CSAI2017P --- Applied Machine Learning (Theory)}\\[2pt]
{\large\bfseries Theory Quiz 3 (T03) \quad\textbullet\quad Set E}\\[3pt]
{\small Duration \textbf{40 minutes} \quad\textbullet\quad Maximum marks \textbf{20}
 \quad\textbullet\quad 15 questions on four pages}
\end{center}
\vspace{2pt}

% Column widths differ slightly from LQ01's: at LQ01's widths the row runs
% ~11.6pt past the margin and "Signature" does not fit its own cell.
\noindent\begin{tabular}{@{}p{1.4cm}p{5.0cm}p{1.9cm}p{3.5cm}p{1.3cm}p{2.4cm}@{}}
\textbf{Name} & \hrulefill & \textbf{SAP ID} & \hrulefill & \textbf{Batch} & \hrulefill\\[9pt]
\textbf{Date} & \hrulefill & \textbf{Signature} & \hrulefill & \textbf{Set} & \fbox{\rule{0pt}{11pt}\ E\ }\\
\end{tabular}

\vspace{4pt}
\noindent\fbox{\parbox{\dimexpr\linewidth-2\fboxsep-2\fboxrule}{\small
\textbf{Syllabus.} Prerequisite material --- probability, Python and elementary calculus --- together with
Units I--III (Lectures 1--3) as retained material: generalisation and data leakage, loss against metric, and
batch sizes and momentum; and, principally, \textbf{Unit III (Lecture 4)}: adaptive optimisers --- Adagrad,
RMSprop, Adadelta and Adam, bias correction, and what adaptivity does and does not buy.\\[3pt]
\textbf{Instructions.} Closed book. Write your answers on this paper, in the space provided. A
calculator is not needed --- the arithmetic is deliberately short. Answer every question --- there
is no negative marking. In Part C a reason is \textbf{optional}, but a correct reason can rescue a
wrong answer. In Parts D and E the marks are for the reasoning, not the vocabulary; an answer that
states what it cannot determine earns credit, one that asserts without a reason does not.\\[3pt]
\textbf{About the set letter.} Sets are handed out \textbf{at random}. Every set covers the same
topics at the same level, for the same marks, against the same marking scheme.}}
\vspace{3pt}


\parthead{A --- Multiple choice}{5 questions $\times$ 1 = 5 marks. Circle one letter.}

\begin{qA}
  \item \texttt{sum(1 for c in deck if c.suit == 'H') / len(deck)}, for a standard $52$-card \texttt{deck}, evaluates to\\
    (a) 0.0769 \quad (b) 0.5\\
    (c) 0.25 \quad (d) 13
  \item After \texttt{xs = [x for x in range(12) if x \% 4 == 0]}, \texttt{len(xs)} is\\
    (a) 4 \quad (b) 3\\
    (c) 2 \quad (d) 12
  \item In \texttt{for xb, yb in loader:} with \texttt{len(dataset) == 4000} and \texttt{batch\_size=32}, the body runs per epoch\\
    (a) 125 times \quad (b) 4000 times\\
    (c) 32 times \quad (d) once
  \item \texttt{Xs = StandardScaler().fit\_transform(X)} is executed before \texttt{train\_test\_split(Xs, y)}. The reported test score is\\
    (a) unaffected --- scaling is deterministic \quad (b) pessimistic --- scaling discards information\\
    (c) an error --- the shapes will not match \quad (d) optimistic --- the test rows helped fit the scaler
  \item In \texttt{s = s + g**2} versus \texttt{s = gamma*s + (1-gamma)*g**2}, the second line is the update used by\\
    (a) Adagrad \quad (b) plain SGD\\
    (c) momentum \quad (d) RMSprop
\end{qA}

\vspace{4pt}

\parthead{B --- Fill in the blanks}{4 $\times$ 1 = 4 marks. Write the missing word on the rule.}

\begin{qB}

  \item \texttt{f = lambda w: w**3}. Then $f'(w)$ is \BL.

  \item \texttt{g.std()} over mini-batches falls as \texttt{batch\_size} grows, like $1/\BL$.

  \item With \texttt{s += g**2} the denominator \texttt{np.sqrt(s)} only ever grows, so the step decays like $1/\BL$.

  \item With \texttt{m\_hat = m/(1-b1**t)} and \texttt{v\_hat = v/(1-b2**t)}, the size of Adam's first step is exactly \BL.

\end{qB}

\vfill
\begin{center}\small
\begin{tabular}{|>{\centering\arraybackslash}m{1.9cm}|*{5}{>{\centering\arraybackslash}m{1.6cm}|}}
\hline
\rule{0pt}{12pt}\textbf{Question} & A1 & A2 & A3 & A4 & A5\\[2pt]\hline
\rule{0pt}{20pt}\textbf{Answer} & & & & & \\[7pt]\hline
\end{tabular}\\[2pt]
{\footnotesize Copy your circled Part A letters into this grid. The grid is what is marked.}
\end{center}

\newpage

\parthead{C --- True or false}{3 $\times$ 1 = 3 marks. Circle one. A reason is optional, but write one --- it can rescue a wrong answer.}

\begin{qC}

  \item \texttt{model.score(X\_test, y\_test)} coming back \emph{lower} than expected is the signature of data leakage. \hfill \textbf{TRUE} \; / \; \textbf{FALSE}\\[1pt]
    {\small Reason:}\,\blk{140mm}\rules{1}{6.4mm}

  \item With \texttt{lr} tuned for each, \texttt{Adam} must reach a lower final loss than \texttt{SGD(momentum=0.9)}. \hfill \textbf{TRUE} \; / \; \textbf{FALSE}\\[1pt]
    {\small Reason:}\,\blk{140mm}\rules{1}{6.4mm}

  \item On a badly conditioned quadratic, \texttt{SGD(momentum=0.9)} can need fewer iterations than \texttt{Adam}. \hfill \textbf{TRUE} \; / \; \textbf{FALSE}\\[1pt]
    {\small Reason:}\,\blk{140mm}\rules{1}{6.4mm}

\end{qC}

\vspace{4pt}

\parthead{D --- Short answer}{2 questions $\times$ 2 = 4 marks. Show how you decide, not just what you decide.}

\begin{qD}

  \item The two loops below are each run for exactly one step from \texttt{s = 0.0}, with \texttt{eta = 0.1} and a first gradient of \texttt{g = -5.0}; \texttt{gamma = 0.96} and \texttt{eps} is negligible.
\begin{code}
# Adagrad                      # RMSprop
s = s + g**2                   s = gamma*s + (1 - gamma)*g**2
w = w - eta*g/np.sqrt(s)       w = w - eta*g/np.sqrt(s)
\end{code}
    \begin{enumerate}[label=(\alph*),leftmargin=7mm,itemsep=1pt,topsep=2pt]
      \item Compute the size of each optimiser's first step, $\lvert \Delta w_1 \rvert$. Give two decimal places where the answer is not exact.
      \item Neither answer depends on how large $g_1$ was. Show why.
    \end{enumerate}\par\vspace{64mm}

\end{qD}

\newpage

\parthead{D --- Short answer, continued}{}

\begin{qD}[start=2]

  \item The training loop minimises \texttt{nn.CrossEntropyLoss()}, but the number in the report comes from \texttt{accuracy\_score(...)}. \textbf{(a)} Explain why these are two different quantities rather than two names for one. \textbf{(b)} State which of them you use to decide the model is finished, and why.\par\vspace{80mm}

\end{qD}

\vspace{4pt}

\parthead{E --- Reasoning}{1 question $\times$ 4 = 4 marks. You are marked on the reasoning you make visible.}

\begin{qE}

  \item A pull request changes the project default from \texttt{SGD(momentum=0.9)} to \texttt{Adam}, with the
justification \emph{``Adam always beats plain gradient descent with momentum''}. You have one afternoon and a
laptop. Design the smallest experiment that would actually settle it.


    \vspace{2mm}
    Answer all four parts overleaf. There is no single right experiment --- you are marked on whether
    the one you describe could actually decide the question.
    \begin{enumerate}[label=(\roman*),leftmargin=7mm,itemsep=2pt,topsep=2pt]
      \item What would you run it on? Say how many problems and of what kind. \hfill \textbf{[1]}
      \item What would you measure, and at what point? \hfill \textbf{[1]}
      \item Name the one thing that must be held equal between the two optimisers, without which the
            comparison means nothing at all. \hfill \textbf{[1]}
      \item State the result that would make you accept the claim, \emph{and} the result that would
            make you reject it. \hfill \textbf{[1]}
    \end{enumerate}

\end{qE}

\newpage

\parthead{E --- Reasoning, answer space}{Label each part (i), (ii), (iii). A sketch, a table or a list is as acceptable as prose.}

\vfill

{\footnotesize\itshape Continue on a continuation sheet if you need more room; label anything you write there with its question number.}

\newpage

\renewcommand{\SetLetter}{F}\setcounter{page}{1}


\begin{center}
{\large\bfseries CSAI2017P --- Applied Machine Learning (Theory)}\\[2pt]
{\large\bfseries Theory Quiz 3 (T03) \quad\textbullet\quad Set F}\\[3pt]
{\small Duration \textbf{40 minutes} \quad\textbullet\quad Maximum marks \textbf{20}
 \quad\textbullet\quad 15 questions on four pages}
\end{center}
\vspace{2pt}

% Column widths differ slightly from LQ01's: at LQ01's widths the row runs
% ~11.6pt past the margin and "Signature" does not fit its own cell.
\noindent\begin{tabular}{@{}p{1.4cm}p{5.0cm}p{1.9cm}p{3.5cm}p{1.3cm}p{2.4cm}@{}}
\textbf{Name} & \hrulefill & \textbf{SAP ID} & \hrulefill & \textbf{Batch} & \hrulefill\\[9pt]
\textbf{Date} & \hrulefill & \textbf{Signature} & \hrulefill & \textbf{Set} & \fbox{\rule{0pt}{11pt}\ F\ }\\
\end{tabular}

\vspace{4pt}
\noindent\fbox{\parbox{\dimexpr\linewidth-2\fboxsep-2\fboxrule}{\small
\textbf{Syllabus.} Prerequisite material --- probability, Python and elementary calculus --- together with
Units I--III (Lectures 1--3) as retained material: generalisation and data leakage, loss against metric, and
batch sizes and momentum; and, principally, \textbf{Unit III (Lecture 4)}: adaptive optimisers --- Adagrad,
RMSprop, Adadelta and Adam, bias correction, and what adaptivity does and does not buy.\\[3pt]
\textbf{Instructions.} Closed book. Write your answers on this paper, in the space provided. A
calculator is not needed --- the arithmetic is deliberately short. Answer every question --- there
is no negative marking. In Part C a reason is \textbf{optional}, but a correct reason can rescue a
wrong answer. In Parts D and E the marks are for the reasoning, not the vocabulary; an answer that
states what it cannot determine earns credit, one that asserts without a reason does not.\\[3pt]
\textbf{About the set letter.} Sets are handed out \textbf{at random}. Every set covers the same
topics at the same level, for the same marks, against the same marking scheme.}}
\vspace{3pt}


\parthead{A --- Multiple choice}{5 questions $\times$ 1 = 5 marks. Circle one letter.}

\begin{qA}
  \item One card is taken at random from an ordinary pack. The chance that it is a heart is\\
    (a) one in thirteen \quad (b) one in two\\
    (c) one in fifty-two \quad (d) one in four
  \item Counting the whole numbers from one to nineteen that leave the same remainder as one when divided by four, the count is\\
    (a) 4 \quad (b) 6\\
    (c) 5 \quad (d) none of them do
  \item Four thousand training examples are worked through thirty-two at a time. In one complete pass, the number of times the weights change is\\
    (a) four thousand \quad (b) one hundred and twenty-five\\
    (c) thirty-two \quad (d) once
  \item Something that learns from the data is allowed to see all of it before any of it is set aside for testing. The score that comes out afterwards is\\
    (a) better than it should be \quad (b) worse than it should be\\
    (c) unchanged \quad (d) impossible to compute
  \item Of the four adaptive methods, the one that asks for no step size at all, working one out for itself from the sizes of the steps it has already taken, is\\
    (a) Adadelta \quad (b) Adagrad\\
    (c) RMSprop \quad (d) Adam
\end{qA}

\vspace{4pt}

\parthead{B --- Fill in the blanks}{4 $\times$ 1 = 4 marks. Write the missing word on the rule.}

\begin{qB}

  \item If a quantity grows as the cube of a number, its rate of change with that number is \BL.

  \item Averaging a gradient over a larger handful of examples makes it less noisy --- the noise falls like one over \BL.

  \item A method whose running total only ever grows must take ever smaller steps: its step falls like one over \BL.

  \item Correcting both running averages for having started at zero makes the very first step exactly \BL\ in size, whatever the gradient was.

\end{qB}

\vfill
\begin{center}\small
\begin{tabular}{|>{\centering\arraybackslash}m{1.9cm}|*{5}{>{\centering\arraybackslash}m{1.6cm}|}}
\hline
\rule{0pt}{12pt}\textbf{Question} & A1 & A2 & A3 & A4 & A5\\[2pt]\hline
\rule{0pt}{20pt}\textbf{Answer} & & & & & \\[7pt]\hline
\end{tabular}\\[2pt]
{\footnotesize Copy your circled Part A letters into this grid. The grid is what is marked.}
\end{center}

\newpage

\parthead{C --- True or false}{3 $\times$ 1 = 3 marks. Circle one. A reason is optional, but write one --- it can rescue a wrong answer.}

\begin{qC}

  \item When something that should not have been seen is seen during training, the tell-tale sign is a disappointing score on the data held back. \hfill \textbf{TRUE} \; / \; \textbf{FALSE}\\[1pt]
    {\small Reason:}\,\blk{140mm}\rules{1}{6.4mm}

  \item Once the step size has been chosen carefully for each method in turn, the ones that adapt per coordinate are bound to come out ahead. \hfill \textbf{TRUE} \; / \; \textbf{FALSE}\\[1pt]
    {\small Reason:}\,\blk{140mm}\rules{1}{6.4mm}

  \item On a long narrow valley, carrying momentum can get you to the bottom in fewer steps than a method that adapts each coordinate separately. \hfill \textbf{TRUE} \; / \; \textbf{FALSE}\\[1pt]
    {\small Reason:}\,\blk{140mm}\rules{1}{6.4mm}

\end{qC}

\vspace{4pt}

\parthead{D --- Short answer}{2 questions $\times$ 2 = 4 marks. Show how you decide, not just what you decide.}

\begin{qD}

  \item Two methods take their very first step from a standing start, both with a step size of $0.03$ and a first gradient of $+8$. The first keeps a running \emph{total} of squared gradients; the second keeps a running \emph{average} of them, forgetting at a rate that leaves one hundredth of the weight on the newest value. Both then divide the step by the square root of what they have accumulated; the guard constant is too small to matter.
    \begin{enumerate}[label=(\alph*),leftmargin=7mm,itemsep=1pt,topsep=2pt]
      \item Compute the size of each optimiser's first step, $\lvert \Delta w_1 \rvert$. Give two decimal places where the answer is not exact.
      \item Neither answer depends on how large $g_1$ was. Show why.
    \end{enumerate}\par\vspace{64mm}

\end{qD}

\newpage

\parthead{D --- Short answer, continued}{}

\begin{qD}[start=2]

  \item A spam filter is trained by making one quantity as small as possible, and is then described to the client by a completely different number --- the fraction of messages it gets right. \textbf{(a)} Explain why these are two different quantities rather than two names for one. \textbf{(b)} State which of them you use to decide the model is finished, and why.\par\vspace{80mm}

\end{qD}

\vspace{4pt}

\parthead{E --- Reasoning}{1 question $\times$ 4 = 4 marks. You are marked on the reasoning you make visible.}

\begin{qE}

  \item Someone claims, flatly, that the method which adapts each coordinate separately always beats the one
that merely carries momentum. You have one afternoon and a laptop. Design the smallest experiment that would
actually settle it.


    \vspace{2mm}
    Answer all four parts overleaf. There is no single right experiment --- you are marked on whether
    the one you describe could actually decide the question.
    \begin{enumerate}[label=(\roman*),leftmargin=7mm,itemsep=2pt,topsep=2pt]
      \item What would you run it on? Say how many problems and of what kind. \hfill \textbf{[1]}
      \item What would you measure, and at what point? \hfill \textbf{[1]}
      \item Name the one thing that must be held equal between the two optimisers, without which the
            comparison means nothing at all. \hfill \textbf{[1]}
      \item State the result that would make you accept the claim, \emph{and} the result that would
            make you reject it. \hfill \textbf{[1]}
    \end{enumerate}

\end{qE}

\newpage

\parthead{E --- Reasoning, answer space}{Label each part (i), (ii), (iii). A sketch, a table or a list is as acceptable as prose.}

\vfill

{\footnotesize\itshape Continue on a continuation sheet if you need more room; label anything you write there with its question number.}

\newpage

\renewcommand{\SetLetter}{G}\setcounter{page}{1}


\begin{center}
{\large\bfseries CSAI2017P --- Applied Machine Learning (Theory)}\\[2pt]
{\large\bfseries Theory Quiz 3 (T03) \quad\textbullet\quad Set G}\\[3pt]
{\small Duration \textbf{40 minutes} \quad\textbullet\quad Maximum marks \textbf{20}
 \quad\textbullet\quad 15 questions on four pages}
\end{center}
\vspace{2pt}

% Column widths differ slightly from LQ01's: at LQ01's widths the row runs
% ~11.6pt past the margin and "Signature" does not fit its own cell.
\noindent\begin{tabular}{@{}p{1.4cm}p{5.0cm}p{1.9cm}p{3.5cm}p{1.3cm}p{2.4cm}@{}}
\textbf{Name} & \hrulefill & \textbf{SAP ID} & \hrulefill & \textbf{Batch} & \hrulefill\\[9pt]
\textbf{Date} & \hrulefill & \textbf{Signature} & \hrulefill & \textbf{Set} & \fbox{\rule{0pt}{11pt}\ G\ }\\
\end{tabular}

\vspace{4pt}
\noindent\fbox{\parbox{\dimexpr\linewidth-2\fboxsep-2\fboxrule}{\small
\textbf{Syllabus.} Prerequisite material --- probability, Python and elementary calculus --- together with
Units I--III (Lectures 1--3) as retained material: generalisation and data leakage, loss against metric, and
batch sizes and momentum; and, principally, \textbf{Unit III (Lecture 4)}: adaptive optimisers --- Adagrad,
RMSprop, Adadelta and Adam, bias correction, and what adaptivity does and does not buy.\\[3pt]
\textbf{Instructions.} Closed book. Write your answers on this paper, in the space provided. A
calculator is not needed --- the arithmetic is deliberately short. Answer every question --- there
is no negative marking. In Part C a reason is \textbf{optional}, but a correct reason can rescue a
wrong answer. In Parts D and E the marks are for the reasoning, not the vocabulary; an answer that
states what it cannot determine earns credit, one that asserts without a reason does not.\\[3pt]
\textbf{About the set letter.} Sets are handed out \textbf{at random}. Every set covers the same
topics at the same level, for the same marks, against the same marking scheme.}}
\vspace{3pt}


\parthead{A --- Multiple choice}{5 questions $\times$ 1 = 5 marks. Circle one letter.}

\begin{qA}
  \item Compare, for one card drawn from a standard pack, $P(\text{a heart})$ with $P(\text{an ace})$\\
    (a) the first is larger \quad (b) the second is larger\\
    (c) they are equal \quad (d) cannot be decided
  \item Compare \texttt{len(range(1, 20, 4))} with \texttt{len('adam')}\\
    (a) the second is larger \quad (b) they are equal\\
    (c) one raises an error \quad (d) the first is larger
  \item Compare the updates per epoch made by $B = 1$ with those made by $B = 32$, on the same data\\
    (a) $B=32$ makes $32\times$ as many \quad (b) they are equal\\
    (c) $B=1$ makes $32\times$ as many \quad (d) it depends on $\eta$
  \item Compare fitting a scaler \emph{inside} each fold with fitting it once on everything beforehand\\
    (a) the second deflates the reported score \quad (b) the second inflates the reported score\\
    (c) they give the same score \quad (d) the second is faster and otherwise identical
  \item Compare what Adagrad accumulates with what RMSprop accumulates\\
    (a) an average, against a sum \quad (b) a sum, against an average\\
    (c) both sums, at different rates \quad (d) both averages, at different rates
\end{qA}

\vspace{4pt}

\parthead{B --- Fill in the blanks}{4 $\times$ 1 = 4 marks. Write the missing word on the rule.}

\begin{qB}

  \item Compare $w^3$ with its derivative: the second is \BL.

  \item Compare the gradient noise at batch size $B$ with that at batch size $1$: the first is smaller by a factor of \BL.

  \item Compare Adagrad's step at step $t$ with its step at step $1$: it has shrunk like $1/\BL$.

  \item Compare Adam's corrected first step with its uncorrected one: the corrected version has size exactly \BL.

\end{qB}

\vfill
\begin{center}\small
\begin{tabular}{|>{\centering\arraybackslash}m{1.9cm}|*{5}{>{\centering\arraybackslash}m{1.6cm}|}}
\hline
\rule{0pt}{12pt}\textbf{Question} & A1 & A2 & A3 & A4 & A5\\[2pt]\hline
\rule{0pt}{20pt}\textbf{Answer} & & & & & \\[7pt]\hline
\end{tabular}\\[2pt]
{\footnotesize Copy your circled Part A letters into this grid. The grid is what is marked.}
\end{center}

\newpage

\parthead{C --- True or false}{3 $\times$ 1 = 3 marks. Circle one. A reason is optional, but write one --- it can rescue a wrong answer.}

\begin{qC}

  \item Between a model with leakage and one without, it is the leaking model that reports the \emph{lower} held-out score. \hfill \textbf{TRUE} \; / \; \textbf{FALSE}\\[1pt]
    {\small Reason:}\,\blk{140mm}\rules{1}{6.4mm}

  \item Between an adaptive optimiser and momentum, each at its own tuned learning rate, the adaptive one always finishes lower. \hfill \textbf{TRUE} \; / \; \textbf{FALSE}\\[1pt]
    {\small Reason:}\,\blk{140mm}\rules{1}{6.4mm}

  \item Between momentum and Adam on a ravine, it can be momentum that needs fewer iterations. \hfill \textbf{TRUE} \; / \; \textbf{FALSE}\\[1pt]
    {\small Reason:}\,\blk{140mm}\rules{1}{6.4mm}

\end{qC}

\vspace{4pt}

\parthead{D --- Short answer}{2 questions $\times$ 2 = 4 marks. Show how you decide, not just what you decide.}

\begin{qD}

  \item Compare the \emph{first} step taken by Adagrad with the first step taken by RMSprop, both from $s_0 = 0$, both with $\eta = 0.6$, and both at a point where the gradient is $g_1 = -3$. Adagrad accumulates $s_t = s_{t-1} + g_t^2$; RMSprop accumulates $s_t = \gamma s_{t-1} + (1-\gamma)g_t^2$ with $\gamma = 0.8$. Both apply $w \leftarrow w - \eta g_t/\sqrt{s_t}$; take $\varepsilon$ as negligible.
    \begin{enumerate}[label=(\alph*),leftmargin=7mm,itemsep=1pt,topsep=2pt]
      \item Compute the size of each optimiser's first step, $\lvert \Delta w_1 \rvert$. Give two decimal places where the answer is not exact.
      \item Neither answer depends on how large $g_1$ was. Show why.
    \end{enumerate}\par\vspace{64mm}

\end{qD}

\newpage

\parthead{D --- Short answer, continued}{}

\begin{qD}[start=2]

  \item Compare the two numbers attached to one classifier: the cross-entropy its training loop minimised, and the accuracy quoted to the client. \textbf{(a)} Explain why these are two different quantities rather than two names for one. \textbf{(b)} State which of them you use to decide the model is finished, and why.\par\vspace{80mm}

\end{qD}

\vspace{4pt}

\parthead{E --- Reasoning}{1 question $\times$ 4 = 4 marks. You are marked on the reasoning you make visible.}

\begin{qE}

  \item Two colleagues disagree. One says \emph{``Adam always beats plain gradient descent with momentum''};
the other says it depends. You have one afternoon and a laptop. Design the smallest experiment that would
actually settle it between them.


    \vspace{2mm}
    Answer all four parts overleaf. There is no single right experiment --- you are marked on whether
    the one you describe could actually decide the question.
    \begin{enumerate}[label=(\roman*),leftmargin=7mm,itemsep=2pt,topsep=2pt]
      \item What would you run it on? Say how many problems and of what kind. \hfill \textbf{[1]}
      \item What would you measure, and at what point? \hfill \textbf{[1]}
      \item Name the one thing that must be held equal between the two optimisers, without which the
            comparison means nothing at all. \hfill \textbf{[1]}
      \item State the result that would make you accept the claim, \emph{and} the result that would
            make you reject it. \hfill \textbf{[1]}
    \end{enumerate}

\end{qE}

\newpage

\parthead{E --- Reasoning, answer space}{Label each part (i), (ii), (iii). A sketch, a table or a list is as acceptable as prose.}

\vfill

{\footnotesize\itshape Continue on a continuation sheet if you need more room; label anything you write there with its question number.}

\newpage

\renewcommand{\SetLetter}{H}\setcounter{page}{1}


\begin{center}
{\large\bfseries CSAI2017P --- Applied Machine Learning (Theory)}\\[2pt]
{\large\bfseries Theory Quiz 3 (T03) \quad\textbullet\quad Set H}\\[3pt]
{\small Duration \textbf{40 minutes} \quad\textbullet\quad Maximum marks \textbf{20}
 \quad\textbullet\quad 15 questions on four pages}
\end{center}
\vspace{2pt}

% Column widths differ slightly from LQ01's: at LQ01's widths the row runs
% ~11.6pt past the margin and "Signature" does not fit its own cell.
\noindent\begin{tabular}{@{}p{1.4cm}p{5.0cm}p{1.9cm}p{3.5cm}p{1.3cm}p{2.4cm}@{}}
\textbf{Name} & \hrulefill & \textbf{SAP ID} & \hrulefill & \textbf{Batch} & \hrulefill\\[9pt]
\textbf{Date} & \hrulefill & \textbf{Signature} & \hrulefill & \textbf{Set} & \fbox{\rule{0pt}{11pt}\ H\ }\\
\end{tabular}

\vspace{4pt}
\noindent\fbox{\parbox{\dimexpr\linewidth-2\fboxsep-2\fboxrule}{\small
\textbf{Syllabus.} Prerequisite material --- probability, Python and elementary calculus --- together with
Units I--III (Lectures 1--3) as retained material: generalisation and data leakage, loss against metric, and
batch sizes and momentum; and, principally, \textbf{Unit III (Lecture 4)}: adaptive optimisers --- Adagrad,
RMSprop, Adadelta and Adam, bias correction, and what adaptivity does and does not buy.\\[3pt]
\textbf{Instructions.} Closed book. Write your answers on this paper, in the space provided. A
calculator is not needed --- the arithmetic is deliberately short. Answer every question --- there
is no negative marking. In Part C a reason is \textbf{optional}, but a correct reason can rescue a
wrong answer. In Parts D and E the marks are for the reasoning, not the vocabulary; an answer that
states what it cannot determine earns credit, one that asserts without a reason does not.\\[3pt]
\textbf{About the set letter.} Sets are handed out \textbf{at random}. Every set covers the same
topics at the same level, for the same marks, against the same marking scheme.}}
\vspace{3pt}


\parthead{A --- Multiple choice}{5 questions $\times$ 1 = 5 marks. Circle one letter.}

\begin{qA}
  \item \emph{One card in four is a heart, and $13/52 = 1/4$.} Drawing one card at random from a standard pack, $P(\text{a heart})$ is\\
    (a) $1/13$ \quad (b) $1/4$\\
    (c) $1/2$ \quad (d) $13$
  \item \emph{\texttt{len(range(1, 10, 3))} is $3$, counting $1, 4, 7$.} Then \texttt{len(range(1, 20, 4))} is\\
    (a) 5 \quad (b) 4\\
    (c) 6 \quad (d) an error
  \item \emph{With $n = 100$ and $B = 10$, one epoch makes $10$ updates.} With $n = 4000$ and $B = 32$, one epoch makes\\
    (a) 4000 \quad (b) 32\\
    (c) 1 \quad (d) 125
  \item \emph{A scaler fitted inside each fold gives an honest score.} The same scaler fitted on everything before the split gives a score that is\\
    (a) too poor \quad (b) identical\\
    (c) too good \quad (d) undefined
  \item \emph{Adagrad keeps a running sum of squared gradients.} RMSprop keeps a running\\
    (a) sum of squared gradients \quad (b) sum of the gradients themselves\\
    (c) average of squared gradients \quad (d) average of the gradients themselves
\end{qA}

\vspace{4pt}

\parthead{B --- Fill in the blanks}{4 $\times$ 1 = 4 marks. Write the missing word on the rule.}

\begin{qB}

  \item \emph{The derivative of $w^2$ is $2w$.} The derivative of $w^3$ is \BL.

  \item \emph{Going from $B=1$ to $B=4$ halves the gradient noise.} In general the noise falls like $1/\BL$.

  \item \emph{After $100$ steps Adagrad's step has shrunk tenfold.} In general the step decays like $1/\BL$.

  \item \emph{Uncorrected, Adam's first step is $3.16\times$ too large at $\beta_2 = 0.999$.} Corrected, its size is exactly \BL.

\end{qB}

\vfill
\begin{center}\small
\begin{tabular}{|>{\centering\arraybackslash}m{1.9cm}|*{5}{>{\centering\arraybackslash}m{1.6cm}|}}
\hline
\rule{0pt}{12pt}\textbf{Question} & A1 & A2 & A3 & A4 & A5\\[2pt]\hline
\rule{0pt}{20pt}\textbf{Answer} & & & & & \\[7pt]\hline
\end{tabular}\\[2pt]
{\footnotesize Copy your circled Part A letters into this grid. The grid is what is marked.}
\end{center}

\newpage

\parthead{C --- True or false}{3 $\times$ 1 = 3 marks. Circle one. A reason is optional, but write one --- it can rescue a wrong answer.}

\begin{qC}

  \item \emph{An honest pipeline gives a believable held-out score.} A leaking one therefore gives a score that is believably \emph{low}. \hfill \textbf{TRUE} \; / \; \textbf{FALSE}\\[1pt]
    {\small Reason:}\,\blk{140mm}\rules{1}{6.4mm}

  \item \emph{On a sparse problem, Adagrad beat plain SGD by a factor of eight.} Therefore an adaptive optimiser, tuned, beats momentum on any problem. \hfill \textbf{TRUE} \; / \; \textbf{FALSE}\\[1pt]
    {\small Reason:}\,\blk{140mm}\rules{1}{6.4mm}

  \item \emph{Adaptive methods measure gradients, and a ravine is a problem of curvature.} So on a ravine, momentum can beat Adam. \hfill \textbf{TRUE} \; / \; \textbf{FALSE}\\[1pt]
    {\small Reason:}\,\blk{140mm}\rules{1}{6.4mm}

\end{qC}

\vspace{4pt}

\parthead{D --- Short answer}{2 questions $\times$ 2 = 4 marks. Show how you decide, not just what you decide.}

\begin{qD}

  \item \emph{From $s_0 = 0$ with $\eta = 1$ and $g_1 = -2$, Adagrad has $s_1 = 4$, so its first step is $1 \times (-2)/2 = -1$: size $1$, which is $\eta$.} Now take $\eta = 0.05$ and $g_1 = +12$, with RMSprop's $\gamma = 0.95$. Adagrad accumulates $s_t = s_{t-1} + g_t^2$; RMSprop accumulates $s_t = \gamma s_{t-1} + (1-\gamma)g_t^2$. Both apply $w \leftarrow w - \eta g_t/\sqrt{s_t}$; take $\varepsilon$ as negligible.
    \begin{enumerate}[label=(\alph*),leftmargin=7mm,itemsep=1pt,topsep=2pt]
      \item Compute the size of each optimiser's first step, $\lvert \Delta w_1 \rvert$. Give two decimal places where the answer is not exact.
      \item Neither answer depends on how large $g_1$ was. Show why.
    \end{enumerate}\par\vspace{64mm}

\end{qD}

\newpage

\parthead{D --- Short answer, continued}{}

\begin{qD}[start=2]

  \item \emph{A ruler and a thermometer both produce numbers, but they are not two names for one measurement.} A classifier's cross-entropy and its accuracy are likewise two different things. \textbf{(a)} Explain why. \textbf{(b)} State which of them you use to decide the model is finished, and why.\par\vspace{80mm}

\end{qD}

\vspace{4pt}

\parthead{E --- Reasoning}{1 question $\times$ 4 = 4 marks. You are marked on the reasoning you make visible.}

\begin{qE}

  \item \emph{Worked analogue. To settle ``drug A cures faster than drug B'' you would not give A to the fit
patients and B to the frail ones --- you would hold everything but the drug equal, and decide in advance what
difference would count.}

\vspace{2mm}
Now settle this one. A classmate asserts: \emph{``Adam always beats plain gradient descent with momentum.''}
You have one afternoon and a laptop. Design the smallest experiment that would actually settle it.


    \vspace{2mm}
    Answer all four parts overleaf. There is no single right experiment --- you are marked on whether
    the one you describe could actually decide the question.
    \begin{enumerate}[label=(\roman*),leftmargin=7mm,itemsep=2pt,topsep=2pt]
      \item What would you run it on? Say how many problems and of what kind. \hfill \textbf{[1]}
      \item What would you measure, and at what point? \hfill \textbf{[1]}
      \item Name the one thing that must be held equal between the two optimisers, without which the
            comparison means nothing at all. \hfill \textbf{[1]}
      \item State the result that would make you accept the claim, \emph{and} the result that would
            make you reject it. \hfill \textbf{[1]}
    \end{enumerate}

\end{qE}

\newpage

\parthead{E --- Reasoning, answer space}{Label each part (i), (ii), (iii). A sketch, a table or a list is as acceptable as prose.}

\vfill

{\footnotesize\itshape Continue on a continuation sheet if you need more room; label anything you write there with its question number.}

\newpage

\end{document}